\documentclass[10pt]{article}
\usepackage[utf8]{inputenc}
\usepackage[T1]{fontenc}
\usepackage[ngerman]{babel}
\usepackage[a4paper,margin=18mm]{geometry}
\usepackage{helvet}
\usepackage{xcolor}
\usepackage{colortbl}
\usepackage{tabularx}
\usepackage{longtable}
\usepackage{array}
\usepackage{enumitem}
\usepackage{fancyhdr}
\usepackage{titlesec}
\usepackage{graphicx}
\usepackage{needspace}
\renewcommand{\familydefault}{\sfdefault}
\definecolor{BEBlue}{HTML}{0066CC}
\definecolor{BEDark}{HTML}{0B2545}
\definecolor{BEOrange}{HTML}{FF9500}
\definecolor{BELight}{HTML}{F5F7FB}
\definecolor{BEPaper}{HTML}{FFF7ED}
\definecolor{BEText}{HTML}{2C3E50}
\definecolor{BEMuted}{HTML}{6C757D}
\pagestyle{fancy}
\fancyhf{}
\renewcommand{\headrulewidth}{0pt}
\lhead{\textcolor{BEDark}{\textbf{Business English in 4 Wochen}}}
\rhead{\textcolor{BEOrange}{\textbf{BE4W}}}
\cfoot{\textcolor{BEMuted}{\thepage}}
\setlength{\parindent}{0pt}
\setlength{\parskip}{5pt}
\emergencystretch=3em
\setlist[itemize]{topsep=3pt,itemsep=3pt}
\setlist[enumerate]{topsep=3pt,itemsep=3pt}
\arrayrulecolor{BEBlue!20}
\titleformat{\section}{\Large\bfseries\color{BEBlue}}{}{0pt}{}[\vspace{-2mm}\textcolor{BEOrange}{\rule{\linewidth}{1.3pt}}]
\titleformat{\subsection}{\large\bfseries\color{BEDark}}{}{0pt}{}
\newcommand{\doctype}[1]{\colorbox{BEPaper}{\parbox{0.96\linewidth}{\vspace{1mm}\textcolor{BEOrange}{\textbf{\MakeUppercase{#1}}}\vspace{1mm}}}\par\vspace{4mm}}
\newcommand{\btitle}[1]{{\Huge\bfseries\color{BEDark} #1}\par\vspace{2mm}\textcolor{BEOrange}{\rule{38mm}{2pt}}\par\vspace{3mm}}
\newcommand{\bsubtitle}[1]{{\large\color{BEText} #1}\par\vspace{7mm}}
\newcommand{\callout}[1]{\vspace{3mm}\fcolorbox{BEOrange!35}{BEPaper}{\parbox{0.94\linewidth}{\raggedright\vspace{2mm}\textcolor{BEText}{#1}\vspace{2mm}}}\vspace{4mm}}
\newcommand{\brandfooter}{\vfill\textcolor{BEMuted}{\small Wort Wach Verlag · shortcutenglish.de · Ergänzendes Material zum Buch}}
\newcommand{\coverblock}[3]{%
  \noindent\colorbox{BEDark}{\parbox{0.96\linewidth}{\raggedright\vspace{5mm}\textcolor{BEOrange}{\small\bfseries\MakeUppercase{#1}}\par\vspace{2mm}{\Huge\bfseries\textcolor{white}{#2}}\par\vspace{2mm}{\large\textcolor{white!82}{#3}}\vspace{5mm}}}\par\vspace{7mm}%
}
\newcommand{\infocard}[2]{%
  \vspace{2mm}\noindent\fcolorbox{BEBlue!18}{white}{\parbox{0.94\linewidth}{\raggedright\vspace{2.5mm}{\large\bfseries\textcolor{BEBlue}{#1}}\par\vspace{1mm}\textcolor{BEText}{#2}\vspace{2.5mm}}}\par\vspace{2mm}%
}
\newcommand{\accentcard}[2]{%
  \vspace{2mm}\noindent\fcolorbox{BEOrange!32}{BEPaper}{\parbox{0.94\linewidth}{\raggedright\vspace{2.5mm}{\large\bfseries\textcolor{BEDark}{#1}}\par\vspace{1mm}\textcolor{BEText}{#2}\vspace{2.5mm}}}\par\vspace{2mm}%
}
\newcommand{\phraseitem}[2]{%
  \noindent\fcolorbox{BEBlue!12}{white}{\parbox{0.94\linewidth}{\raggedright\vspace{1.8mm}\textcolor{BEMuted}{\small Deutsch}\par\textcolor{BEText}{#1}\par\vspace{1.2mm}\textcolor{BEOrange}{\small Englisch}\par{\bfseries\textcolor{BEDark}{#2}}\vspace{1.8mm}}}\par\vspace{2mm}%
}
\newcommand{\safechapter}[1]{\Needspace{10\baselineskip}\section{#1}}
\newcommand{\writingcard}[2]{%
  \vspace{2mm}\noindent\fcolorbox{BEOrange!32}{BEPaper}{\parbox{0.94\linewidth}{\raggedright\vspace{2.5mm}{\large\bfseries\textcolor{BEDark}{#1}}\par\vspace{#2}}}\par\vspace{2mm}%
}
\newcolumntype{Y}{>{\raggedright\arraybackslash}X}

\begin{document}
\coverblock{Woche 1--4 · Abruf-Workbook}{Core-Phrasen Workbook}{172 Business-English-Phrasen zum Abdecken, Abrufen und Wiederholen}
\callout{So trainierst du mit dem Kartenlayout: Lies zuerst nur die deutsche Aufgabe. Decke die englische Lösung mit einem Zettel, deiner Hand oder einem zweiten Blatt ab. Formuliere die englische Phrase laut aus dem Kopf. Erst danach prüfst du die Lösung und markierst schwierige Karten.}
\accentcard{Warum Kartenlayout?}{Dieses Workbook ist nicht zum passiven Durchlesen gedacht. Es folgt dem Coach-Prinzip: Deutsch sehen, Englisch abrufen, Lösung prüfen. Für eigene Listen, Anki, Excel oder Notion sind die CSV-Dateien besser geeignet.}
\safechapter{Woche 1 · Sicher starten · 45 Phrasen}
\phraseitem{Bevor wir weitermachen, ein kurzer Punkt.}{Before we move on, one quick point.}
\phraseitem{Könnten Sie bitte erläutern, was Sie mit \ldots{} meinen?}{Could you clarify what you mean by \ldots{}?}
\phraseitem{Von meiner Seite aus ist der nächste Schritt \ldots{}}{From my side, the next step is \ldots{}}
\phraseitem{Können wir fortfahren?}{Are we good to proceed?}
\phraseitem{Nur um sicherzugehen, dass wir auf einer Linie sind \ldots{}}{Just to make sure we're aligned \ldots{}}
\phraseitem{Lassen Sie mich kurz zusammenfassen.}{Let me quickly summarize.}
\phraseitem{Ich schicke nach dem Gespräch eine kurze Zusammenfassung.}{I'll share a quick summary after the call.}
\phraseitem{Könnten wir uns \ldots{} genauer ansehen?}{Could we take a closer look at \ldots{}?}
\phraseitem{Sind wir uns einig?}{Are we on the same page?}
\phraseitem{Lassen Sie uns mit dem nächsten Punkt weitermachen.}{Let's move on to the next topic.}
\phraseitem{Danke für die Aktualisierung.}{Thanks for the update.}
\phraseitem{Um es zusammenzufassen \ldots{}}{So to recap \ldots{}}
\phraseitem{Aus meiner Sicht \ldots{}}{From my perspective \ldots{}}
\phraseitem{Das passt für mich.}{That works for me.}
\phraseitem{Klingt das für Sie gut?}{Does that sound good to you?}
\phraseitem{Lassen Sie uns das später heute abstimmen.}{Let's align on this later today.}
\phraseitem{Ich melde mich morgen dazu.}{I'll follow up tomorrow.}
\phraseitem{Könnten Sie mir die aktuelle Version schicken?}{Could you send me the latest version?}
\phraseitem{Lassen Sie uns einen kurzen Call vereinbaren.}{Let's set up a quick call.}
\phraseitem{Nur um sicherzugehen \ldots{}}{Just to double-check \ldots{}}
\phraseitem{Spricht etwas dagegen, wenn wir weitermachen?}{Any objections if we continue?}
\phraseitem{Wir liegen etwas hinter dem Zeitplan.}{We're slightly behind schedule.}
\phraseitem{Wir liegen im Plan für Freitag.}{We're on track for Friday.}
\phraseitem{Lassen Sie es kurz und klar halten.}{Let's keep it short and clear.}
\phraseitem{Das ist ein guter Punkt.}{That's a good point.}
\phraseitem{Lassen Sie mich etwas ergänzen.}{Let me add something to that.}
\phraseitem{Ich stimme vollkommen zu.}{I totally agree.}
\phraseitem{Ich verstehe, was Sie meinen.}{I see what you mean.}
\phraseitem{Könnten Sie ein kurzes Beispiel geben?}{Could you give a quick example?}
\phraseitem{Lassen Sie uns auf die wichtigsten Punkte konzentrieren.}{Let's focus on the key points.}
\phraseitem{Das ist eine interessante Idee.}{That's an interesting idea.}
\phraseitem{Lassen Sie das nächste Meeting planen.}{Let's schedule the next meeting.}
\phraseitem{Wir finalisieren die Präsentation bis Freitag.}{We'll finalize the deck by Friday.}
\phraseitem{Bitte teilen Sie Ihren Bildschirm.}{Please share your screen.}
\phraseitem{Können mich alle hören?}{Can everyone hear me?}
\phraseitem{Entschuldigung, du bist stummgeschaltet.}{Sorry, you're on mute.}
\phraseitem{Lassen Sie uns mit einem kurzen Update beginnen.}{Let's start with a quick update.}
\phraseitem{Wer möchte beginnen?}{Who wants to start?}
\phraseitem{Lassen Sie uns reihum gehen.}{Let's go around the table.}
\phraseitem{Von meiner Seite ist das klar.}{That's clear from my side.}
\phraseitem{Gibt es weitere Anmerkungen?}{Any other comments?}
\phraseitem{Lassen Sie uns hier abschließen.}{Let's wrap up here.}
\phraseitem{Der nächste Schritt ist \ldots{}}{The next step is \ldots{}}
\phraseitem{Danke an alle, gute Runde.}{Thanks everyone, good session.}
\phraseitem{Ich freue mich auf das nächste Meeting.}{Looking forward to the next one.}
\safechapter{Woche 2 · Professioneller klingen · 58 Phrasen}
\phraseitem{Danke für das Update.}{Thanks for the update.}
\phraseitem{Könnten wir die Dateien eventuell morgen teilen?}{Could we possibly share the files tomorrow?}
\phraseitem{Das würde uns sehr helfen, am Freitag zu starten.}{That would really help us start on Friday.}
\phraseitem{Zwei kurze Punkte und ein Vorschlag.}{Two quick points and a proposal.}
\phraseitem{Also, um es zusammenzufassen \ldots{}}{So, just to recap \ldots{}}
\phraseitem{Lassen wir das für später.}{Let's park this for later.}
\phraseitem{Wäre es machbar, \ldots{}?}{Would it be feasible to \ldots{}?}
\phraseitem{Würden Sie Ihre Gedanken dazu teilen?}{Would you mind sharing your thoughts?}
\phraseitem{Das könnte schwierig werden.}{That might be tricky.}
\phraseitem{Das sollten wir uns vielleicht noch einmal ansehen.}{We might want to revisit that.}
\phraseitem{Könnten wir kurz einen Schritt zurückgehen?}{Could we take a step back for a second?}
\phraseitem{Ich wollte nur kurz nachfragen \ldots{}}{I just wanted to double-check \ldots{}}
\phraseitem{Wenn ich etwas hinzufügen darf \ldots{}}{If I may add something \ldots{}}
\phraseitem{Würde es Sinn machen, \ldots{}?}{Would it make sense to \ldots{}?}
\phraseitem{Das klingt für mich vernünftig.}{That sounds reasonable to me.}
\phraseitem{Ich verstehe Ihren Punkt vollkommen.}{I completely see your point.}
\phraseitem{Das ist ein berechtigter Punkt.}{That's a fair point.}
\phraseitem{Wir schweifen ein wenig ab.}{We're getting slightly off track.}
\phraseitem{Könnten wir später darauf zurückkommen?}{Could we come back to that later?}
\phraseitem{Lassen Sie uns zum nächsten Punkt übergehen.}{Let's move to the next item.}
\phraseitem{Nur um es klarzustellen \ldots{}}{Just to clarify \ldots{}}
\phraseitem{Wenn ich Sie richtig verstanden habe \ldots{}}{If I understood you correctly \ldots{}}
\phraseitem{Es wäre vielleicht überlegenswert \ldots{}}{It might be worth considering \ldots{}}
\phraseitem{Das ist eine interessante Sichtweise.}{That's an interesting perspective.}
\phraseitem{Ich nehme das auf.}{I'll take that on board.}
\phraseitem{Lassen Sie uns akzeptieren, dass wir uneins sind.}{Let's agree to disagree.}
\phraseitem{Da sind wir uns einig.}{We're aligned on that.}
\phraseitem{Ich verstehe, wie Sie das meinen.}{I see where you're coming from.}
\phraseitem{Könnten Sie das näher erläutern?}{Could you elaborate on that?}
\phraseitem{Würden Sie sagen, das hat Priorität?}{Would you say that's a priority?}
\phraseitem{Ich würde vorschlagen, wir halten es kurz.}{I'd suggest we keep it short.}
\phraseitem{Lassen Sie uns beim Hauptthema bleiben.}{Let's stick to the main topic.}
\phraseitem{Nur um die Erwartungen zu steuern \ldots{}}{Just to manage expectations \ldots{}}
\phraseitem{Das ist mir nicht ganz klar.}{That's not entirely clear to me.}
\phraseitem{Könnten Sie das vielleicht umformulieren?}{Could you perhaps rephrase that?}
\phraseitem{Das ist im Moment vielleicht nicht ideal.}{That might not be ideal right now.}
\phraseitem{Lassen Sie uns das außerhalb der Runde besprechen.}{Let's take this offline.}
\phraseitem{Danke, das ist hilfreich.}{Thanks, that's helpful.}
\phraseitem{Ich sehe kurz in meinen Notizen nach.}{Let me quickly check my notes.}
\phraseitem{Ich schätze Ihre Geduld.}{I appreciate your patience.}
\phraseitem{Das ist sehr freundlich von Ihnen.}{That's very kind of you.}
\phraseitem{Ich befürchte, das wird nicht funktionieren.}{I'm afraid that won't work.}
\phraseitem{Das ist ein berechtigtes Anliegen.}{That's a valid concern.}
\phraseitem{Könnten wir das auf nächste Woche verschieben?}{Could we postpone this to next week?}
\phraseitem{Schauen wir, was wir tun können.}{Let's see what we can do.}
\phraseitem{Ich bin mir da nicht ganz sicher.}{I'm not entirely sure about that.}
\phraseitem{Das ist auf jeden Fall einen Blick wert.}{That's definitely worth exploring.}
\phraseitem{Lassen Sie uns das im Hinterkopf behalten.}{Let's keep that in mind.}
\phraseitem{Wir kommen später darauf zurück.}{We'll come back to that later.}
\phraseitem{Das können wir uns näher ansehen.}{That's something we can look into.}
\phraseitem{Klingt das vernünftig?}{Does that sound reasonable?}
\phraseitem{Wäre das für alle in Ordnung?}{Would that work for everyone?}
\phraseitem{Ich denke, wir sind uns im Großen und Ganzen einig.}{I think we're broadly in agreement.}
\phraseitem{Das ist ein konstruktiver Vorschlag.}{That's a constructive suggestion.}
\phraseitem{Ich notiere das.}{I'll make a note of that.}
\phraseitem{Könnten Sie mir die Details anschließend schicken?}{Could you send me the details afterwards?}
\phraseitem{Ich würde Ihr Feedback dazu schätzen.}{I'd appreciate your feedback on this.}
\phraseitem{Ziel ist es, das bis Freitag abzuschließen.}{Let's aim to finalize this by Friday.}
\safechapter{Woche 3 · Flexibel reagieren · 49 Phrasen}
\phraseitem{Können wir das bis Freitag festmachen?}{Can we lock this by Friday?}
\phraseitem{Lassen Sie uns eine Lösung finden.}{Let's make it work.}
\phraseitem{Ich komme später heute darauf zurück.}{I'll circle back later today.}
\phraseitem{Ich verstehe Ihren Punkt, aber \ldots{}}{I see your point, but \ldots{}}
\phraseitem{Lassen Sie uns das aus einer anderen Perspektive betrachten.}{Let's take a different angle.}
\phraseitem{Guter Punkt -- hier ist, was ich ergänzen würde.}{Good point -- here's what I'd add.}
\phraseitem{Wir machen gute Fortschritte.}{We're making solid progress.}
\phraseitem{Lassen Sie uns dieses Momentum nutzen.}{Let's keep that momentum.}
\phraseitem{Lassen Sie uns kurz ehrlich sein.}{Let's be real for a second.}
\phraseitem{Können wir das beschleunigen?}{Can we speed this up?}
\phraseitem{Lassen Sie uns direkt zur Sache kommen.}{Let's get straight to the point.}
\phraseitem{Was ist der schnellste Weg, das zu lösen?}{What's the quickest way to solve this?}
\phraseitem{Lassen Sie uns aktiv werden.}{Let's take action.}
\phraseitem{Lassen Sie uns das schnell angehen.}{Let's move fast on this.}
\phraseitem{Können wir uns darauf festlegen?}{Can we commit to that?}
\phraseitem{Lassen Sie uns alle ins Boot holen.}{Let's get everyone on board.}
\phraseitem{Hier ist mein Vorschlag.}{Here's what I propose.}
\phraseitem{Lassen Sie uns nicht zu viel darüber nachdenken.}{Let's not overthink this.}
\phraseitem{Was blockiert uns?}{What's blocking us?}
\phraseitem{Lassen Sie uns das heute lösen.}{Let's fix this today.}
\phraseitem{Wir brauchen eine klare Entscheidung.}{We need a clear decision.}
\phraseitem{Was ist hier unser größtes Risiko?}{What's our biggest risk here?}
\phraseitem{Lassen Sie uns das schnell abstimmen.}{Let's align quickly.}
\phraseitem{Ich übernehme das.}{I can take ownership of that.}
\phraseitem{Wir müssen das voranbringen.}{We need to move this forward.}
\phraseitem{Lassen Sie uns das aufteilen.}{Let's break this down.}
\phraseitem{Können wir ein kurzes Update bekommen?}{Can we get a quick update?}
\phraseitem{Das ist kein Ausschlusskriterium.}{That's not a dealbreaker.}
\phraseitem{Lassen Sie uns konzentriert bleiben.}{Let's stay focused.}
\phraseitem{Lassen Sie uns den Gang wechseln}{Let's shift gears.}
\phraseitem{Ich bin zuversichtlich, dass wir das schaffen.}{I'm confident we can do this.}
\phraseitem{Wo stehen wir?}{Where do we stand?}
\phraseitem{Lassen Sie uns das heute abschließen.}{Let's close this today.}
\phraseitem{Mir gefällt die Richtung.}{I like where this is going.}
\phraseitem{Lassen Sie uns das ins Ziel bringen.}{Let's push this across the finish line.}
\phraseitem{Lassen Sie uns loslegen.}{Let's move on this.}
\phraseitem{Wir sind in einer guten Position.}{We're in a good spot.}
\phraseitem{Das hat hohe Priorität.}{This is high priority.}
\phraseitem{Lassen Sie uns die nächsten Schritte klären.}{Let's clarify next steps.}
\phraseitem{Perfekt -- lassen Sie uns das so machen.}{Perfect -- let's go with that.}
\phraseitem{Ich kümmere mich darum.}{I'm on it.}
\phraseitem{Klingt gut für mich.}{Sounds good to me.}
\phraseitem{Lassen Sie uns hier die Führung übernehmen.}{Let's take the lead here.}
\phraseitem{Wie sieht unser Zeitplan aus?}{What's our timeline?}
\phraseitem{Wir brauchen eine schnelle Entscheidung.}{We need a quick decision.}
\phraseitem{Lassen Sie uns den Plan noch einmal ansehen.}{Let's revisit the plan.}
\phraseitem{Wir eskalieren das, falls nötig.}{Let's escalate if needed.}
\phraseitem{Lassen Sie uns später heute abstimmen.}{Let's sync later today.}
\phraseitem{Wir brauchen Klarheit, bevor wir weitermachen.}{We need clarity before moving on.}
\safechapter{Woche 4 · Im Meeting anwenden · 20 Phrasen}
\phraseitem{Kurze Notiz zur Bestätigung \ldots{}}{Quick note to confirm \ldots{}}
\phraseitem{Ich schätze Ihr Feedback.}{Appreciate your feedback.}
\phraseitem{Ich nehme Sie hier mit ins Boot.}{Looping you in here.}
\phraseitem{Nur um sicherzugehen, dass wir abgestimmt sind.}{Just to make sure we're aligned.}
\phraseitem{Die wichtigste Erkenntnis ist \ldots{}}{Main takeaway is \ldots{}}
\phraseitem{Der nächste Schritt von meiner Seite \ldots{}}{Next step from my side \ldots{}}
\phraseitem{Ich freue mich auf Ihre Rückmeldung.}{Looking forward to your thoughts.}
\phraseitem{Nochmals danke für Ihre Zeit.}{Thanks again for your time.}
\phraseitem{Hier ein kurzes Update dazu.}{Here's a quick update on this.}
\phraseitem{Lassen Sie mich einen Punkt klarstellen.}{Let me clarify one point.}
\phraseitem{Ich teile hier die aktuelle Version.}{Sharing the latest version here.}
\phraseitem{Anbei die Datei.}{Attaching the file below.}
\phraseitem{Geben Sie Bescheid, wenn etwas unklar ist.}{Let me know if anything's unclear.}
\phraseitem{Ich beantworte gerne alle Fragen.}{Happy to answer any questions.}
\phraseitem{Ich schätze die schnelle Rückmeldung.}{Appreciate the quick turnaround.}
\phraseitem{Das fällt mir besonders auf.}{Here's what stands out to me.}
\phraseitem{Ich bestätige kurz unseren Plan für Freitag.}{Just confirming our plan for Friday.}
\phraseitem{Kommentieren Sie gerne.}{Feel free to add comments.}
\phraseitem{Anschluss an unser Gespräch von vorhin.}{Following up on our call earlier.}
\phraseitem{Ich teile bald weitere Details.}{I'll share more details soon.}
\brandfooter
\end{document}
